\section{Numerical Simulation of Non-Ideal Components}
\label{sec:nonideal}

\subsection{Nonlinear Modeling and State-Space Representation}

The analytic solutions derived in the previous sections rely on ideal,
constant component values.
In practical circuits, however, passive components can change behavior under
high-stress transient conditions.
To test how these material limits alter the ideal mathematical model, this
section adds two phenomenological nonlinear effects: current-dependent
heating in the resistor and magnetic saturation in the inductor.

Joule heating causes a resistor's temperature, and therefore its effective
resistance, to increase as current flows through it.
We model this dynamic resistance as
\begin{equation}
  R(i) = R_0 + k i^2,
  \label{eq:thermal_resistance}
\end{equation}
where $R_0$ is the nominal resistance and $k$ is a heating coefficient.
Similarly, an inductor with a saturating magnetic core loses effective
inductance at high current.
This effect is approximated by
\begin{equation}
  L(i) = \frac{L_0}{1 + (i/I_{\mathrm{sat}})^2},
  \label{eq:saturating_inductance}
\end{equation}
where $L_0$ is the nominal inductance and $I_{\mathrm{sat}}$ is the saturation
current threshold.

Substituting these dynamic parameters into Kirchhoff's Voltage Law gives the
effective nonlinear model
\begin{equation}
  L(i)C\frac{d^2\vC}{dt^2} + R(i)C\frac{d\vC}{dt} + \vC = V(t),
  \label{eq:nonlinear_rlc}
\end{equation}
with series current $i(t)=C\,d\vC/dt$.
Because the resistive term now contains a cubic dependence on
$d\vC/dt$ and the inductive term changes with current, the characteristic
equation method used for constant-coefficient systems no longer applies.

For numerical simulation, let $y_1=\vC$ and $y_2=d\vC/dt$.
Then $i=C y_2$, and \eqref{eq:nonlinear_rlc} becomes the first-order
state-space system
\begin{align}
  \frac{dy_1}{dt} &= y_2, \label{eq:nonideal_state_1}\\
  \frac{dy_2}{dt} &=
    \frac{V(t)-y_1-R(Cy_2)Cy_2}{L(Cy_2)C}.
    \label{eq:nonideal_state_2}
\end{align}
The system was integrated in Python using an adaptive RK45 method with
$R_0=50\,\Omega$, $L_0=0.1\,\mathrm{H}$, $C=10\,\mu\mathrm{F}$,
$k=8000$, and $I_{\mathrm{sat}}=0.05\,\mathrm{A}$.
These values are not meant to identify a specific commercial component;
instead, they isolate how stress-dependent component behavior changes the
same baseline RLC circuit studied analytically above.

\begin{figure}[!t]
  \centering
  \includegraphics[width=\textwidth]{nonideal_rlc_simulation}
  \caption{Numerical step responses of the RLC circuit under non-ideal
           material constraints compared against the ideal baseline model
           ($R_0=50\,\Omega$, $L_0=0.1\,\mathrm{H}$,
           $C=10\,\mu\mathrm{F}$).}
  \label{fig:nonideal_simulation}
\end{figure}

\subsection{Simulation Results and Physical Interpretation}

The numerical simulations in Fig.~\ref{fig:nonideal_simulation} show that
the ideal step response remains a useful baseline, but it does not capture
all of the transient behavior produced by stressed real components.

\subsubsection{Thermal Resistor Response}

During the initial transient interval, the initially uncharged capacitor draws
a large inrush current.
In the thermal model, this current surge raises the effective resistance above
the nominal $50\,\Omega$ value exactly when the circuit has the most stored
transient energy.
The resistor therefore acts like a dynamic brake: it dissipates energy more
aggressively during the first swing, suppressing the overshoot from about
$14.44\,\mathrm{V}$ in the ideal model to about $13.07\,\mathrm{V}$ and
bringing the response back toward the $10\,\mathrm{V}$ steady state more
quickly.

\subsubsection{Inductor Saturation Response}

When inductor saturation is introduced, the response also deviates from the
ideal baseline, but for a different reason.
The high initial current lowers the effective inductance in
\eqref{eq:saturating_inductance}; since
$\omega_0=1/\sqrt{LC}$, the instantaneous natural frequency increases while
the inductor is saturated.
The voltage peak therefore occurs earlier than in the ideal model, and the
first overshoot is reduced to about $13.24\,\mathrm{V}$.
As the transient current decays, $L(i)$ returns toward $L_0$, so the response
gradually realigns with the ideal curve.

\subsubsection{AC Input at Resonance}

To demonstrate energy buildup, the step input was replaced with an AC source
\begin{equation}
  V(t) = \Vs \sin(\omega_0 t),
  \label{eq:resonant_ac_input}
\end{equation}
driven at the ideal natural frequency $\omega_0=1000\,\mathrm{rad/s}$.
Because the driving frequency matches the circuit's natural energy exchange
between $L$ and $C$, the capacitor voltage grows far beyond the source
amplitude before settling into a large oscillation.
In this simulation, the capacitor voltage approaches nearly
$\pm20\,\mathrm{V}$ even though the source amplitude is only $10\,\mathrm{V}$.

These simulations refine the research question.
An ideal resistor determines the broad damping regime, but real-world
component nonlinearities make damping state-dependent: the resistor, inductor,
and input waveform can all change the transient path taken toward stability.
